\scanpage{183}
On prend
\[
 \varphi(f,g)=\widetilde C(f\,dg)
\]

a) $\Lambda$-bilinéaire : clair

b) nul si \deleted{$f$ section de $\Lambda$} $f=1$ ou $g=1$ : clair, car $\widetilde C$ est nul sur les diff. exactes

c) $\varphi$ est antisymétrique, i.e. $\left\{\begin{array}{l}\text{alterné si }p\ne2\\\text{symétrique si }p=2\end{array}\right.$

i.e. $\varphi(f,g)+\varphi(g,f)=0$ i.e. $\widetilde C(f\,dg+g\,df)=0$, en effet $f\,dg+g\,df=d(fg)=0$.

\uline{Proposition} L'accouplement $(\ast)$ est une dualité parfaite de $\Lambda$-modules.

On peut supposer qu'il existe une $p$-base relative $\{x\}$, et on a une base de $Q$ sur $1$ formé par les classes des \deleted{$x,x^{\unsure{2|3}},\ldots,x^{p-1}$}. $x^i$ $1\le i\le p-1$. On a
\[
 \widetilde C(x^i\,dx^j)=\widetilde C(jx^{i+j-1}\,dx)=
 \begin{cases}0&\text{si }i+j\ne p-1\\j&\text{si }i+j=p-1\end{cases}
\]
donc la matrice de la forme $\varphi$ est
\[
 \begin{pmatrix}
 &0&&p-1\\
 &&\reflectbox{$\ddots$}&\\
 &2&&0\\
 1&&&
 \end{pmatrix},
\]
son déterminant, au signe près, est donc $1.2.\ldots(p-1)\in\mathbb F_p^\ast$, \margin{$\Delta=-1\in\mathbb F_p^\ast$} donc inversible, O.K.

\uline{Corollaire} \deleted{$\det Q$} On a un isom canonique
\[
 Q\xrightarrow{\sim}\check Q\otimes(\Omega^1_{X^{p/S}/S}\otimes\Lambda)
\]
$\Downarrow$

d'où un isom canonique
\[
 \det Q\simeq\det\check Q\otimes(\Omega^1_{X^{p/S}/S}{}^{\otimes(p-1)}\otimes\Lambda)
\]
% archive sheet number: 182
